% GENERATED FILE. Edit canonical lessons, then run npm run generate:books.
% canonical-source-hash: fnv1a64:9b6928bf679320a3
% canonical-lessons: FR-C512-chorale, FR-C512-fredonner, FR-C512-siffloter, FR-C512-refrain, FR-C512-melodie

\chapter{Chanter juste}
\label{ch:fr-chanter-juste}
\hlchaptermodality{512}

\begin{chapteropening}
\textbf{By the end of this chapter:} \emph{I can talk about singing and tunes.}

Complete the last lesson of chapter 512: I can talk about singing and tunes.
\end{chapteropening}

\section[la chorale]{la chorale — choir}

\label{lesson:FR-C512-chorale}



\begin{quote}
Before the new one: say the French for accordion, then the French for to tune.
\end{quote}

\subsection*{You'll want to know: la chorale}
\textbf{la chorale} — \textquotedblleft{}choir\textquotedblright{}.

\begin{quote}
Je chante dans une chorale depuis trois ans.
\end{quote}

\begin{quote}
\emph{I have been singing in a choir for three years.}
\end{quote}

An amateur choir; \emph{le chœur} is a choir too, also in a church.

\begin{cousinweb}
\emph{la chorale} comes from \emph{choral}, built on Medieval Latin \emph{chorālis}, from Latin \emph{chorus} \textquotedblleft{}a dance, a choir\textquotedblright{}, a borrowing from Greek \emph{khorós}.
\end{cousinweb}

\begin{grammarlens}[title={What you've built}]
One more word for telling, explaining and describing at B1.
\end{grammarlens}

\subsection*{Your turn}
\begin{itemize}
\raggedright
  \item \emph{Say it:} \emph{la chorale}
  \item \emph{Say it:} \emph{la chorale}, once more
  \item \emph{Say it:} a sentence of your own with \emph{la chorale}
  \item \emph{Recall:} say the French for accordion, then the French for to tune, then say \emph{la chorale} again
\end{itemize}

\begin{tcolorbox}[breakable,colback=teal!4,colframe=teal!35!black,title={Before you move on}]
What does la chorale mean? (\textquotedblleft{}Choir\textquotedblright{}.) Say it once more.
\end{tcolorbox}
\section[fredonner]{fredonner — to hum}

\label{lesson:FR-C512-fredonner}



\begin{quote}
Before the new one: say the French for to tune, then the French for choir.
\end{quote}

\subsection*{You'll want to know: fredonner}
\textbf{fredonner} — \textquotedblleft{}to hum\textquotedblright{}.

\begin{quote}
Elle fredonne une chanson en faisant la cuisine.
\end{quote}

\begin{quote}
\emph{She hums a song while she cooks.}
\end{quote}

Humming with the lips closed or almost closed.

\begin{cousinweb}
\emph{fredonner} is built on old \emph{fredon} \textquotedblleft{}a trill\textquotedblright{} (EtymDB), which is usually traced to Latin \emph{fritinnīre} \textquotedblleft{}to twitter\textquotedblright{} (Lewis \& Short); the origin is not fully settled.
\end{cousinweb}

\begin{grammarlens}[title={What you've built}]
One more word for telling, explaining and describing at B1.
\end{grammarlens}

\subsection*{Your turn}
\begin{itemize}
\raggedright
  \item \emph{Say it:} \emph{fredonner}
  \item \emph{Say it:} \emph{fredonner}, once more
  \item \emph{Say it:} a sentence of your own with \emph{fredonner}
  \item \emph{Recall:} say the French for to tune, then the French for choir, then say \emph{fredonner} again
\end{itemize}

\begin{tcolorbox}[breakable,colback=teal!4,colframe=teal!35!black,title={Before you move on}]
What does fredonner mean? (\textquotedblleft{}To hum\textquotedblright{}.) Say it once more.
\end{tcolorbox}
\section[siffloter]{siffloter — to whistle softly (a tune)}

\label{lesson:FR-C512-siffloter}



\begin{quote}
Before the new one: say the French for choir, then the French for to hum.
\end{quote}

\subsection*{You'll want to know: siffloter}
\textbf{siffloter} — \textquotedblleft{}to whistle softly (a tune)\textquotedblright{}.

\begin{quote}
Mon grand-père sifflote en marchant.
\end{quote}

\begin{quote}
\emph{My grandfather whistles softly as he walks.}
\end{quote}

Built on \emph{siffler} with \emph{-oter}, which makes the action light and repeated.

\begin{cousinweb}
\emph{siffloter} is built on \emph{siffler} \textquotedblleft{}to whistle\textquotedblright{} with \emph{-oter}; \emph{siffler} goes back to Latin \emph{sībilāre} \textquotedblleft{}to hiss, to whistle\textquotedblright{} (Lewis \& Short).
\end{cousinweb}

\begin{grammarlens}[title={What you've built}]
One more word for telling, explaining and describing at B1.
\end{grammarlens}

\subsection*{Your turn}
\begin{itemize}
\raggedright
  \item \emph{Say it:} \emph{siffloter}
  \item \emph{Say it:} \emph{siffloter}, once more
  \item \emph{Say it:} a sentence of your own with \emph{siffloter}
  \item \emph{Recall:} say the French for choir, then the French for to hum, then say \emph{siffloter} again
\end{itemize}

\begin{tcolorbox}[breakable,colback=teal!4,colframe=teal!35!black,title={Before you move on}]
What does siffloter mean? (\textquotedblleft{}To whistle softly (a tune)\textquotedblright{}.) Say it once more.
\end{tcolorbox}
\section[le refrain]{le refrain — chorus, refrain}

\label{lesson:FR-C512-refrain}



\begin{quote}
Before the new one: say the French for to hum, then the French for to whistle softly.
\end{quote}

\subsection*{You'll want to know: le refrain}
\textbf{le refrain} — \textquotedblleft{}chorus, refrain\textquotedblright{}.

\begin{quote}
Tout le monde chante le refrain avec lui.
\end{quote}

\begin{quote}
\emph{Everyone sings the chorus with him.}
\end{quote}

The verses are \emph{les couplets}. \emph{C'est toujours le même refrain} — it is always the same old story.

\begin{cousinweb}
\emph{le refrain} comes from Old French \emph{refraindre} \textquotedblleft{}to break off, to repeat\textquotedblright{} (EtymDB), from Vulgar Latin \emph{refrangere}, built on \emph{frangere} \textquotedblleft{}to break\textquotedblright{} (Lewis \& Short).
\end{cousinweb}

\begin{grammarlens}[title={What you've built}]
One more word for telling, explaining and describing at B1.
\end{grammarlens}

\subsection*{Your turn}
\begin{itemize}
\raggedright
  \item \emph{Say it:} \emph{le refrain}
  \item \emph{Say it:} \emph{le refrain}, once more
  \item \emph{Say it:} a sentence of your own with \emph{le refrain}
  \item \emph{Recall:} say the French for to hum, then the French for to whistle softly, then say \emph{le refrain} again
\end{itemize}

\begin{tcolorbox}[breakable,colback=teal!4,colframe=teal!35!black,title={Before you move on}]
What does le refrain mean? (\textquotedblleft{}Chorus, refrain\textquotedblright{}.) Say it once more.
\end{tcolorbox}
\section[la mélodie]{la mélodie — melody, tune}

\label{lesson:FR-C512-melodie}



\begin{quote}
Before the new one: say the French for to whistle softly, then the French for chorus.
\end{quote}

\subsection*{You'll want to know: la mélodie}
\textbf{la mélodie} — \textquotedblleft{}melody, tune\textquotedblright{}.

\begin{quote}
Cette mélodie me rappelle mon enfance.
\end{quote}

\begin{quote}
\emph{This tune reminds me of my childhood.}
\end{quote}

The adjective is \emph{mélodieux} \textquotedblleft{}tuneful\textquotedblright{}.

\begin{cousinweb}
\emph{la mélodie} comes from Late Latin \emph{melōdia}, a borrowing from Greek \emph{melōidía} \textquotedblleft{}singing\textquotedblright{}, built on \emph{mélos} \textquotedblleft{}song\textquotedblright{} and \emph{ōid\'{\={e}}} \textquotedblleft{}ode\textquotedblright{}.
\end{cousinweb}

\begin{grammarlens}[title={What you've built}]
One more word for telling, explaining and describing at B1.
\end{grammarlens}

\subsection*{Your turn}
\begin{itemize}
\raggedright
  \item \emph{Say it:} \emph{la mélodie}
  \item \emph{Say it:} \emph{la mélodie}, once more
  \item \emph{Say it:} a sentence of your own with \emph{la mélodie}
  \item \emph{Recall:} say the French for to whistle softly, then the French for chorus, then say \emph{la mélodie} again
\end{itemize}

\begin{tcolorbox}[breakable,colback=teal!4,colframe=teal!35!black,title={Before you move on}]
What does la mélodie mean? (\textquotedblleft{}Melody, tune\textquotedblright{}.) Say it once more.
\end{tcolorbox}
